\documentclass[conference]{IEEEtran}
\IEEEoverridecommandlockouts
\usepackage{cite}
\usepackage{amsmath,amssymb,amsfonts}
\usepackage{algorithmic}
\usepackage{graphicx}
\usepackage{textcomp}
\usepackage{xcolor}
\usepackage{booktabs}
\usepackage{multirow}
\usepackage{array}
\usepackage{url}
\def\BibTeX{{\rm B\kern-.05em{\sc i\kern-.025em b}\kern-.08em
    T\kern-.1667em\lower.7ex\hbox{E}\kern-.125emX}}
\begin{document}

\title{Double-Gated Time-Window Aggregation: A Multimodal AI Approach to Classroom Attention Scoring}

\author{\IEEEauthorblockN{1\textsuperscript{st} B.V. Tejas}
\IEEEauthorblockA{\textit{Dept. of Computer Science and Engineering} \\
\textit{PES University}\\
Bengaluru, India}
\and
\IEEEauthorblockN{2\textsuperscript{nd} Krishna Keerthan Reddy}
\IEEEauthorblockA{\textit{Dept. of Computer Science and Engineering} \\
\textit{PES University}\\
Bengaluru, India}
\and
\IEEEauthorblockN{3\textsuperscript{rd} R. Saurav Srinivas}
\IEEEauthorblockA{\textit{Dept. of Computer Science and Engineering} \\
\textit{PES University}\\
Bengaluru, India}
}

\maketitle

% ============================================================
% ABSTRACT (strictly grounded in project goals and actual outputs)
% ============================================================
\begin{abstract}
Single-modality systems for automated classroom engagement tracking suffer from well-documented blind spots. Camera-based approaches miss verbal participation entirely, while audio-only methods cannot distinguish on-topic reasoning from irrelevant chatter. We present a dual-stream multimodal pipeline that addresses both failure modes. The visual stream extracts five 3D skeletal landmarks via MediaPipe Pose and classifies physical posture through geometric heuristics, supplemented by YOLOv8-nano object detection for cell phone and beverage identification. The audio stream chains OpenAI Whisper transcription, a Semi-Supervised Contrastive (SSCon) dialogue act classifier built on DialoGPT-medium contextual embeddings and XGBoost, and a Sentence-Transformer cosine similarity gate against the instructor's lesson text. Both streams are aggregated into synchronized one-minute time windows and fused via a weighted 60/40 blend. We report results from controlled classroom sessions where the system correctly differentiates a student providing on-topic reasoning (scoring 88/100) from a student producing off-topic speech with identical physical posture (scoring 48/100).
\end{abstract}

\begin{IEEEkeywords}
Multimodal Learning Analytics, Pose Estimation, Dialogue Act Classification, Contrastive Learning, Semantic Similarity, Attention Scoring
\end{IEEEkeywords}

% ============================================================
% I. INTRODUCTION
% ============================================================
\section{Introduction}

Measuring student engagement during live instruction remains a persistent challenge for educators and educational technology researchers alike. Manual observation does not scale beyond small groups, and the subjective nature of human judgment introduces inconsistency across observers \cite{dewan2019}.
% [INSERT REAL CITATION FOR: survey on engagement detection limitations]

Automated approaches have largely followed one of two paths. Vision-based systems use cameras to track head pose, gaze direction, or body posture and infer attentiveness from physical cues \cite{zaletelj2017}.
% [INSERT REAL CITATION FOR: body pose for classroom attention prediction]
These systems perform adequately for detecting gross distractions---a student turned completely away from the board, for instance---but they are blind to verbal engagement. A student can maintain perfect posture while contributing nothing to the discussion.

Audio-based systems take the opposite approach: they capture speech and reward vocal participation \cite{suresh2022}.
% [INSERT REAL CITATION FOR: TalkMoves or dialogue act classification in classrooms]
The weakness here is equally clear. A student asking their neighbor about weekend plans generates speech activity that these systems cannot distinguish from genuine academic reasoning.

Neither modality alone tells the full story. We built a system that watches and listens simultaneously, and cross-validates one stream against the other within discrete time windows. The audio pathway does not simply count utterances. Every piece of speech must pass through two sequential gates: an acoustic classifier that identifies reasoning-type dialogue acts, and a semantic similarity check that confirms the content relates to the lesson being taught. This double-gating mechanism is the primary defense against students gaming the participation metric.

The contributions of this work are:
\begin{itemize}
    \item A geometry-based posture classifier using five MediaPipe Pose 3D landmarks (nose, left/right ears, left/right shoulders) that replaces a prior CNN approach which exhibited strong bias toward the ``listening'' class.
    \item A double-gated NLP pipeline that chains SSCon dialogue act classification with Sentence-Transformer cosine similarity scoring against instructor-provided lesson text.
    \item A time-window aggregation and fusion framework that synchronizes visual and audio data into one-minute epochs and produces a single weighted engagement score per student.
    \item A proportional behavioral feedback generator that produces natural-language teacher dashboard labels based on per-student dialogue act distributions.
\end{itemize}

% ============================================================
% II. RELATED WORK
% ============================================================
\section{Related Work}

\subsection{Vision-Based Engagement Detection}
Camera-based student monitoring has been explored through facial action units, eye tracking, and full-body pose estimation. Zaletelj and Ko\v{s}ir used OpenPose body features combined with gaze estimation for attention prediction in recorded classroom videos \cite{zaletelj2017}.
% [INSERT REAL CITATION FOR: Zaletelj & Kosir, EURASIP 2017, body pose attention prediction]
Dewan et al.\ provided a broad survey of engagement detection methods, noting that most visual approaches operate on a single modality and do not incorporate speech content \cite{dewan2019}.
% [INSERT REAL CITATION FOR: Dewan et al., Smart Learning Environments 2019, engagement detection review]
More recent work has adopted lightweight models such as YOLOv8 for real-time person and object detection in educational settings \cite{jocher2023}.
% [INSERT REAL CITATION FOR: Ultralytics YOLOv8, 2023]

\subsection{Dialogue Act Classification in Education}
Suresh et al.\ proposed the TalkMoves framework, which classified teacher and student utterances into pedagogically meaningful dialogue acts using fine-tuned transformer models \cite{suresh2022}.
% [INSERT REAL CITATION FOR: Suresh et al., AAAI 2022, TalkMoves]
Their results demonstrated that acts such as \emph{Providing Evidence} and \emph{Making a Claim} can be reliably identified from classroom transcripts, but they did not fuse this information with any visual modality.

\subsection{Sentence Embeddings and Semantic Similarity}
Reimers and Gurevych introduced Sentence-BERT, showing that Siamese BERT-based architectures produce embeddings suitable for efficient cosine similarity comparison \cite{reimers2019}.
% [INSERT REAL CITATION FOR: Reimers & Gurevych, EMNLP 2019, Sentence-BERT]
The \texttt{all-MiniLM-L6-v2} variant maps sentences to a 384-dimensional vector space and has become a standard baseline for semantic textual similarity tasks.

\subsection{Pose Estimation Frameworks}
Lugaresi et al.\ introduced MediaPipe, a modular framework for perception pipelines that includes a Pose Landmarker providing 33 3D landmarks per detected person \cite{lugaresi2019}.
% [INSERT REAL CITATION FOR: Lugaresi et al., arXiv 2019, MediaPipe]
The Pose Landmarker Lite variant used in this work runs efficiently on edge devices and consumer laptops.

\subsection{Automatic Speech Recognition}
Radford et al.\ released Whisper, a general-purpose speech recognition model trained on 680{,}000 hours of multilingual web audio \cite{radford2023}.
% [INSERT REAL CITATION FOR: Radford et al., ICML 2023, Whisper]
The base variant (74M parameters) provides a practical balance between transcription quality and inference speed for real-time classroom applications.

To our knowledge, no prior system combines geometric pose classification, dialogue act recognition, and semantic relevance gating within a unified time-window aggregation framework.

% ============================================================
% III. SYSTEM ARCHITECTURE / METHODOLOGY
% ============================================================
\section{System Architecture}

The system runs two parallel processing streams coordinated by a central session runner. The computer vision pipeline executes as a separate subprocess, while the audio pipeline runs in a daemon thread within the main Python process. This process-level isolation is an engineering requirement: the deep learning frameworks used by each stream (the YOLO/MediaPipe stack and the PyTorch-based NLP stack) conflict when loaded into shared memory. A \texttt{threading.Event} object coordinates shutdown across both streams. At session end, a report generation module parses the independent log files produced by each pipeline.

\subsection{Stream 1: Computer Vision Pipeline}

The visual stream identifies students by face and classifies their physical behavior. It operates on every 30th captured frame to limit computational load.

\subsubsection{Person Detection}
Each sampled frame is processed by YOLOv8-nano (6.5\,MB model file) \cite{jocher2023}. The detector produces bounding boxes for every \texttt{person} class instance. It simultaneously flags the presence of two distraction-related classes: \texttt{cell phone} and \texttt{bottle}/\texttt{cup}.

\subsubsection{Geometric Posture Classification}
For each detected person, the bounding box crop is converted from BGR to RGB and passed to a MediaPipe Pose Landmarker Lite model (5.8\,MB task bundle) \cite{lugaresi2019}. Of the 33 returned 3D landmarks, we use exactly five, listed in Table~\ref{tab:landmarks}.

\begin{table}[htbp]
\caption{MediaPipe Landmarks Used for Classification}
\begin{center}
\begin{tabular}{|c|c|l|}
\hline
\textbf{Index} & \textbf{Landmark} & \textbf{Purpose} \\
\hline
0 & Nose & Head position anchor \\
7 & Left Ear & Horizontal turn indicator \\
8 & Right Ear & Horizontal turn indicator \\
11 & Left Shoulder & Vertical baseline (left) \\
12 & Right Shoulder & Vertical baseline (right) \\
\hline
\end{tabular}
\label{tab:landmarks}
\end{center}
\end{table}

Three geometric metrics are derived from these landmarks. The horizontal distances from the nose to each ear:
\begin{equation}
d_L = |x_{\text{nose}} - x_{\text{left\_ear}}|, \quad d_R = |x_{\text{nose}} - x_{\text{right\_ear}}|
\label{eq:ear_dist}
\end{equation}

The vertical drop of the nose below the shoulder midpoint:
\begin{equation}
\Delta_y = \frac{y_{\text{left\_shoulder}} + y_{\text{right\_shoulder}}}{2} - y_{\text{nose}}
\label{eq:nose_drop}
\end{equation}

Classification proceeds in strict priority order. Table~\ref{tab:rules} lists the exact rules as implemented in the codebase.

\begin{table}[htbp]
\caption{Posture Classification Rules (Priority Order)}
\begin{center}
\begin{tabular}{|c|p{4.0cm}|c|c|}
\hline
\textbf{Rule} & \textbf{Condition} & \textbf{Label} & \textbf{Conf.} \\
\hline
R1a & Either ear presence $< 0.3$ & Distracted & 90\% \\
\hline
R1b & $\frac{\max(d_L, d_R)}{\min(d_L, d_R)} > 3.0$ & Distracted & 85\% \\
\hline
R2 & $\Delta_y < 0.15$ & Writing & 90\% \\
\hline
R3 & All above pass (upright) & Listening & $\geq$70\% \\
\hline
R4 & No landmarks detected & Distracted & 50\% \\
\hline
\end{tabular}
\label{tab:rules}
\end{center}
\end{table}

For R3, the confidence is computed as $\min(99.0,\; 70.0 + \Delta_y \times 100)$. In practice, an upright student produces $\Delta_y$ values in the range 0.25--0.40, yielding confidences of 95--99\%.

\subsubsection{YOLO Object Overrides}
After pose classification, YOLO detections override the label when specific objects appear. If a \texttt{cell phone} is detected and the current label is \texttt{listen} or \texttt{write}, it is forced to \texttt{distracted} at 95\% confidence. If a \texttt{bottle} or \texttt{cup} is detected instead, the label becomes \texttt{drink} at 95\%. Phone detection takes strict priority over beverage detection.

\subsubsection{Face Recognition}
Each person crop is also processed by a face recognition module that computes 128-dimensional face encodings and matches them against pre-registered student photographs using Euclidean distance. Results are logged to a CSV file with columns: \texttt{Timestamp}, \texttt{Person}, \texttt{Behavior}, \texttt{Confidence}.

\subsection{Stream 2: Audio and NLP Pipeline}

The audio stream captures ambient classroom speech, transcribes it, classifies the dialogue act, and checks semantic relevance against the lesson text.

\subsubsection{Speech Capture and Transcription}
Audio is recorded in 10-second chunks at 16{,}000\,Hz (160{,}000 samples per chunk) via the \texttt{sounddevice} library. Each chunk is saved as a temporary WAV file and transcribed by OpenAI Whisper (base model, 74M parameters) \cite{radford2023}. The transcription call sets \texttt{condition\_on\_previous\_text=False} to suppress hallucinated repetitive output during silence. Utterances shorter than three characters are discarded.

\subsubsection{Speaker Attribution}
The system attributes each utterance to the student most recently visible in the camera stream by reading the last row of the CV log CSV file and extracting the \texttt{Person} column value.

\subsubsection{SSCon Dialogue Act Classification}
Each valid transcript is classified into one of four dialogue acts using a Semi-Supervised Contrastive (SSCon) pipeline. The four act categories are listed in Table~\ref{tab:acts}.

\begin{table}[htbp]
\caption{SSCon Dialogue Act Categories}
\begin{center}
\begin{tabular}{|c|l|c|}
\hline
\textbf{Index} & \textbf{Dialogue Act} & \textbf{Is Reasoning} \\
\hline
0 & Relating to another student & Yes \\
1 & Asking for more information & No \\
2 & Making a claim & Yes \\
3 & Providing evidence & Yes \\
\hline
\end{tabular}
\label{tab:acts}
\end{center}
\end{table}

Three out of four acts are treated as ``reasoning.'' Only ``Asking for more information'' is excluded. The classification pipeline proceeds in three stages:

\textbf{Stage 1: Contextual Encoding.} The conversation history (up to 10 prior turns) is concatenated with EOS tokens and fed to DialoGPT-medium \cite{zhang2020}.
% [INSERT REAL CITATION FOR: Zhang et al., ACL 2020, DialoGPT]
A 1024-dimensional hidden-state vector is extracted from layer 12 at the final token position.

\textbf{Stage 2: Sentence Embedding.} The current utterance is independently embedded by a Sentence-Transformer (\texttt{all-MiniLM-L6-v2}) \cite{reimers2019} into a 384-dimensional vector.

\textbf{Stage 3: Projection and Classification.} The two vectors are concatenated into a 1408-dimensional feature vector, projected to 384 dimensions by a trained contrastive bi-encoder, and classified by a 4-class XGBoost model.

The bi-encoder architecture is a six-layer feedforward network: Linear(1408$\rightarrow$720) $\rightarrow$ [BN $\rightarrow$ Linear $\rightarrow$ BN $\rightarrow$ ReLU $\rightarrow$ Dropout(0.2)] $\times$ 4, narrowing progressively from 720 to 640 to 560 to 384 dimensions. Training uses binary cross-entropy on cosine similarity labels, Adam optimizer at $\text{lr} = 7 \times 10^{-4}$, batch size 512, and a distillation probability threshold of 0.85. The XGBoost classifier uses \texttt{multi:softprob} objective with 500 estimators and the \texttt{hist} tree method.

\subsubsection{Semantic Relevance Gate}
After dialogue act classification, a second gate checks whether the spoken content relates to the lesson. Both the transcript and the instructor's lesson text are encoded by the Sentence-Transformer into 384-dimensional vectors. Their cosine similarity is computed:
\begin{equation}
\text{sim}(\mathbf{v}_{\text{stu}}, \mathbf{v}_{\text{les}}) = \frac{\mathbf{v}_{\text{stu}} \cdot \mathbf{v}_{\text{les}}}{\|\mathbf{v}_{\text{stu}}\| \; \|\mathbf{v}_{\text{les}}\|}
\label{eq:cosine}
\end{equation}

An utterance is classified as on-topic if $\text{sim} \geq 0.35$. This threshold was set deliberately low; classroom speech is informal and will rarely match lesson text verbatim. An empty subject text file disables the gate entirely (returns true for all utterances).

% ============================================================
% IV. SCORING MODEL AND AGGREGATION
% ============================================================
\section{Scoring Model and Aggregation}

\subsection{CV Stream: Majority Voting}
Frame-level behavior labels are grouped into one-minute buckets. The number of buckets is $K = \lceil N / 60 \rceil$ where $N$ is the total number of logged frames. For bucket $t$ containing frame labels $b_1, b_2, \ldots, b_n$, the dominant behavior is:
\begin{equation}
B_t = \text{mode}(b_1, b_2, \ldots, b_n)
\label{eq:mode}
\end{equation}

Each dominant behavior maps to a point value:

\begin{table}[htbp]
\caption{CV Behavior Point Values}
\begin{center}
\begin{tabular}{|l|c|l|}
\hline
\textbf{Behavior} & \textbf{Points} & \textbf{Description} \\
\hline
Writing & 5 & Actively taking notes \\
Listening & 4 & Attentive to board/teacher \\
Drinking & 1 & Brief, mild distraction \\
Distracted & 0 & Disengaged \\
\hline
\end{tabular}
\label{tab:cv_points}
\end{center}
\end{table}

Maximum possible score per bucket is 5, yielding $P_{\text{CV,max}} = 5K$.

\subsection{Audio Stream: Tiered Double-Gate Scoring}
Audio events are grouped into one-minute buckets aligned with the CV session length. Within each bucket, every utterance is scored through the double gate according to the tiered rules in Table~\ref{tab:tiers}.

\begin{table}[htbp]
\caption{Tiered Audio Scoring}
\begin{center}
\begin{tabular}{|c|p{3.4cm}|c|}
\hline
\textbf{Tier} & \textbf{Condition} & \textbf{Points} \\
\hline
1 & Reasoning act AND $\text{sim} \geq 0.35$ & 5 \\
\hline
2 & Not reasoning, but $\text{sim} \geq 0.35$ & 3 \\
\hline
3 & Short acknowledgment only & 1 \\
\hline
0 & Off-topic or silence & 0 \\
\hline
\end{tabular}
\label{tab:tiers}
\end{center}
\end{table}

Each bucket retains only the highest tier score among its utterances. If the score reaches 5, evaluation for that bucket terminates early. The acknowledgment set contains 18 tokens: \{yes, no, yeah, yep, nope, ok, okay, sure, right, correct, true, false, hmm, ya, na, haan, nahi\}. Silent buckets score zero. $P_{\text{Audio,max}} = 5K$.

\subsection{Weighted Score Fusion}
Both streams are converted to percentages and blended:
\begin{equation}
C_{\%} = \frac{P_{\text{CV,earned}}}{P_{\text{CV,max}}} \times 100, \quad
A_{\%} = \frac{P_{\text{Audio,earned}}}{P_{\text{Audio,max}}} \times 100
\label{eq:pct}
\end{equation}
\begin{equation}
E = 0.60 \cdot C_{\%} + 0.40 \cdot A_{\%}
\label{eq:blend}
\end{equation}

The 60/40 split reflects that visual data is continuous while audio events are sparse. Fallback logic handles degenerate cases: if CV data is absent, the audio weight becomes 1.0 (and vice versa). If both streams are empty, $E = 0$.

\subsection{Behavioral Feedback Generation}
A feedback generator counts per-student utterances by tier and applies threshold rules. Both absolute counts and proportional checks are used so that the rules remain meaningful for sessions of varying length.

\begin{table}[htbp]
\caption{Feedback Threshold Rules}
\begin{center}
\begin{tabular}{|c|p{3.4cm}|p{2.6cm}|}
\hline
\textbf{\#} & \textbf{Trigger} & \textbf{Label} \\
\hline
1 & No speaking events & Silent Observer \\
\hline
2 & Questioning $> 4$ or $> 40\%$ & Highly Inquisitive \\
\hline
3 & Reasoning $> 3$ or $> 40\%$ & Strong Analytical \\
\hline
4 & Ack.\ $>$ reasoning + quest. & Passive Participant \\
\hline
5 & Off-topic $> 2$ or $> 50\%$ & Distracted Speaker \\
\hline
\end{tabular}
\label{tab:feedback}
\end{center}
\end{table}

Multiple rules may fire simultaneously; their labels are concatenated. If none triggers, the default is ``Standard Participant.''

% ============================================================
% V. EXPERIMENTAL RESULTS & EVALUATION
% ============================================================
\section{Experimental Results}

\subsection{Test Environment}
The system was evaluated in controlled classroom sessions with four enrolled students. % [MISSING DATA: Please insert exact number of sessions conducted]
Hardware was an Apple MacBook Pro with M3 Pro processor running macOS. PyTorch operations used Apple Metal Performance Shaders (MPS) acceleration. Face encodings were pre-registered from a single photograph per student. The lesson topic was Newtonian gravity; the reference text (494 characters) covered gravitational force, mass dependence, and the inverse-square law.

\subsection{Double-Gate Discrimination}
Table~\ref{tab:gate_demo} presents utterances from a single test session demonstrating the semantic gate's ability to separate on-topic reasoning from off-topic speech.

\begin{table}[htbp]
\caption{Semantic Gate: On-Topic vs.\ Off-Topic Discrimination}
\begin{center}
\begin{tabular}{|p{0.9cm}|p{2.1cm}|p{1.5cm}|c|c|}
\hline
\textbf{Student} & \textbf{Utterance} & \textbf{SSCon Act} & \textbf{Sim.} & \textbf{Pts} \\
\hline
Tejas & ``Gravity is important because it pulls objects...'' & Providing evidence & 0.62 & 5 \\
\hline
Krishna & ``When are we going to go for lunch?'' & Asking for info & 0.12 & 0 \\
\hline
Krishna & ``When is the break?'' & Asking for info & 0.09 & 0 \\
\hline
\end{tabular}
\label{tab:gate_demo}
\end{center}
\end{table}

Tejas's statement passes both gates: SSCon labels it as \emph{Providing evidence} (a reasoning act), and the cosine similarity of 0.62 exceeds the 0.35 threshold. Krishna's two utterances both fail the semantic gate---cosine similarities of 0.12 and 0.09 fall well below the threshold---and receive zero audio points despite being valid speech acts. A system without the semantic gate would have credited Krishna for talking.

\subsection{Multimodal Score Fusion}
Table~\ref{tab:fused} shows the fused engagement scores from the same session.

\begin{table}[htbp]
\caption{Fused Engagement Scores (Same Session)}
\begin{center}
\begin{tabular}{|l|c|c|c|l|}
\hline
\textbf{Student} & \textbf{CV} & \textbf{Audio} & \textbf{Final} & \textbf{Feedback} \\
\hline
Tejas & 80\% & 100\% & 88.00 & Analytical \\
\hline
Krishna & 80\% & 0\% & 48.00 & Distracted \\
\hline
Saurav & 100\% & 20\% & 68.00 & Passive \\
\hline
\end{tabular}
\label{tab:fused}
\end{center}
\end{table}

Tejas and Krishna had identical CV percentages (80\%), yet their final scores differ by 40 points. The gap originates entirely from the audio stream. A vision-only system would have scored them identically. Saurav presents a third profile: high physical attentiveness (dominant behavior was \texttt{write}, scoring 100\%) but only brief acknowledgment-level speech (three ``yes/no'' responses), producing the ``Passive Participant'' feedback label.

\subsection{System Output}
The verbatim terminal report for two representative students is reproduced below.

\begin{verbatim}
====================================================
    FINAL STUDENT ATTENTIVENESS REPORT
====================================================

Student: Tejas
Physical Attentiveness (CV)  : 4/5 (80.00%)
  Frames logged: 52
  Dominant state: listen
Verbal Participation (Audio) : 5/5 (100.00%)
  Utterances: 2
  Breakdown: Reasoning(1) | Yes/No(1)
  Feedback: Strong Analytical Skills
FINAL SCORE: 88.00 / 100 (60% CV | 40% Audio)

----------------------------------------------------

Student: krishna
Physical Attentiveness (CV)  : 4/5 (80.00%)
  Frames logged: 45
  Dominant state: listen
Verbal Participation (Audio) : 0/5 (0.00%)
  Utterances: 2
  Breakdown: Questioning(2) | Off-Topic(2)
  Feedback: Distracted Speaker
FINAL SCORE: 48.00 / 100 (60% CV | 40% Audio)
\end{verbatim}

\subsection{YOLO Object Override}
During one test, a student held a cell phone directly in front of the webcam while seated upright. The pose classifier returned \texttt{listen} at 99\%+ confidence---the student's shoulders and nose geometry indicated an attentive posture. With the YOLO override active, the \texttt{cell phone} detection forced the label to \texttt{distracted} at 95\% confidence. The same mechanism was validated for \texttt{bottle} and \texttt{cup} detections, which correctly triggered the \texttt{drink} label.

\subsection{Comparison with Prior CNN Approach}
The original system used a convolutional neural network trained on 112$\times$112 pixel crops to classify posture into four classes. This CNN exhibited strong bias toward the \texttt{listen} class, predicting it at 99\%+ confidence even when the student's head was turned sideways or tilted down toward a phone. Replacing the CNN with the five-point geometric analysis resolved these misclassifications. The geometric approach explicitly encodes head rotation via ear distances and head tilt via nose-shoulder drop---spatial information that pixel-level classification does not reliably capture.
% [MISSING DATA: Please insert CNN training details (epochs, dataset size, architecture layers) if available]

\subsection{Model Inventory}
Table~\ref{tab:models} lists all models deployed in the system with their storage requirements.

\begin{table}[htbp]
\caption{Deployed Model Inventory}
\begin{center}
\begin{tabular}{|l|l|r|}
\hline
\textbf{Component} & \textbf{Model} & \textbf{Size} \\
\hline
Object Detection & YOLOv8-nano & 6.5\,MB \\
Pose Estimation & MediaPipe Pose Lite & 5.8\,MB \\
Speech Recognition & Whisper base (74M) & $\sim$140\,MB \\
Sentence Embedding & all-MiniLM-L6-v2 & $\sim$80\,MB \\
Context Encoding & DialoGPT-medium & $\sim$1.5\,GB \\
Bi-Encoder Projection & SSCon (6-layer FFN) & 10.8\,MB \\
Act Classification & XGBoost (4-class) & 1.5\,MB \\
\hline
\end{tabular}
\label{tab:models}
\end{center}
\end{table}

% ============================================================
% VI. CONCLUSION & FUTURE WORK
% ============================================================
\section{Conclusion and Future Work}

We presented a working multimodal engagement tracking system that fuses geometric pose estimation with double-gated dialogue analysis inside one-minute time windows. The system runs on consumer hardware without dedicated GPU acceleration. Controlled tests show that it correctly differentiates genuine academic engagement from superficial participation---a distinction that neither modality achieves independently.

Several limitations point to directions for future work. First, speaker attribution currently relies on the last student detected by the camera, a simplification that breaks down when multiple students speak in rapid succession; integrating a dedicated speaker diarization model would address this. Second, the cosine similarity threshold of 0.35 was set heuristically for a physics lesson and may require adjustment for subjects with different vocabulary characteristics---automatically calibrating this threshold per subject area is a natural extension. Third, the system operates on a single camera and microphone; scaling to multi-sensor arrays would enable coverage of larger classrooms. Finally, moving from post-session terminal reports to a real-time web dashboard would allow instructors to intervene during class rather than reviewing engagement data after the fact.

% ============================================================
% REFERENCES
% ============================================================
\begin{thebibliography}{00}

\bibitem{jocher2023} G. Jocher, A. Chaurasia, and J. Qiu, ``Ultralytics YOLOv8,'' 2023. [Online]. Available: https://github.com/ultralytics/ultralytics
% [INSERT REAL CITATION FOR: Ultralytics YOLOv8 — use the official GitHub or Zenodo DOI if available]

\bibitem{lugaresi2019} C. Lugaresi et al., ``MediaPipe: A framework for building perception pipelines,'' \emph{arXiv preprint arXiv:1906.08172}, 2019.
% [INSERT REAL CITATION FOR: confirm arXiv ID and add workshop/conference venue if published]

\bibitem{reimers2019} N. Reimers and I. Gurevych, ``Sentence-BERT: Sentence embeddings using Siamese BERT-Networks,'' in \emph{Proc. EMNLP}, 2019, pp. 3982--3992.
% [INSERT REAL CITATION FOR: confirm page numbers and proceedings volume]

\bibitem{suresh2022} A. Suresh et al., ``Using transformers to provide teachers with personalized feedback on their classroom discourse: The TalkMoves application,'' in \emph{Proc. AAAI}, vol. 36, no. 11, 2022, pp. 12828--12835.
% [INSERT REAL CITATION FOR: confirm volume, issue, page numbers]

\bibitem{radford2023} A. Radford, J. W. Kim, T. Xu, G. Brockman, C. McLeavey, and I. Sutskever, ``Robust speech recognition via large-scale weak supervision,'' in \emph{Proc. ICML}, 2023, pp. 28492--28518.
% [INSERT REAL CITATION FOR: confirm ICML 2023 proceedings page numbers]

\bibitem{zhang2020} Y. Zhang et al., ``DialoGPT: Large-scale generative pre-training for conversational response generation,'' in \emph{Proc. ACL System Demonstrations}, 2020, pp. 270--278.
% [INSERT REAL CITATION FOR: confirm ACL 2020 demo track page numbers]

\bibitem{dewan2019} M. A. A. Dewan, M. Murshed, and F. Lin, ``Engagement detection in online learning: A review,'' \emph{Smart Learning Environments}, vol. 6, no. 1, pp. 1--20, 2019.
% [INSERT REAL CITATION FOR: confirm journal volume and DOI]

\bibitem{zaletelj2017} J. Zaletelj and A. Ko\v{s}ir, ``Predicting students' attention in the classroom from Kinect facial and body features,'' \emph{EURASIP J. Image Video Process.}, vol. 2017, no. 1, pp. 1--12, 2017.
% [INSERT REAL CITATION FOR: confirm EURASIP journal volume and DOI]

\end{thebibliography}

\end{document}
